\section{Conclusion}
\label{sec:conclusion}

An enforcement architecture in which an award-layer screening
stage triages firms and environments before a costly bid-layer
forensic stage interrogates them is feasible, and the data
quantify what the architecture buys. On São Paulo's BEC
($\valYearStart$--$\valYearEnd$), routing forensic
interrogation through a frequent-loser flag built on
contract-award records reduces the bid-microdata pool the
forensic stage must work on by $\valArchFootprintReduction$
($\valArchSeqKone$ of $\valArchPoolSize$ firms), while still
recovering $\valArchTPThouSeqK$ of $\valArchNPos$ adjudicated
cobidders. That is the headline, and it is an
enforcement-design claim before it is a screening-technology
claim. The screen is a triage device, not an adjudication
device: it ranks loser-side firms and procurement environments
for costly bid-layer interrogation, not cartel members for
legal sanction. Its discrimination accuracy against the
cobidder population (firm-level AUC $\valAUCFLfirmTemp$ under
temporal holdout, $\valAUCprePost$ on the strict pre-$2020$
benchmark) is the technology that makes the architecture
feasible; its complementarity with bid-distribution screens
($+\valImhofIncFL$ AUC over the seven-feature
Imhof--Wallimann pipeline, DeLong $p=\valImhofIncFLp$) is what
makes the sequencing well-defined. The simple
separating-equilibrium argument we develop motivates endogenous
loser-side participation as the ranking primitive on the award
layer; the binary $\text{median}+1.5\!\times\!\text{IQR}$ rule
on always-losers operationalises it. Cobidders match the
modelled cover-bidder type along four of five firm-level
predictions; at the bid level they bid plausibly close to
winners with elevated within-firm dispersion, the signature of
credible cover bidding rather than the textbook
deliberately-uncompetitive form.

The screening statistic and the bid-distribution moment battery
are informational complements, not substitutes. They
discriminate the same cartel-adjacency target on different
observability layers and carry non-redundant signal when both
are accessed; the screening stage functions as a credible
Stage-1 gatekeeper for the forensic stage. Reforms that mandate
operational bid-microdata archival expand the forensic stage,
not the screening one. Jurisdictions whose operational layer
carries award records but not per-bidder bid amounts can deploy
the screening stage on data they already maintain, which
reverses the standard policy intuition that operational
bid-archival is the rate-limiting precondition for proactive
cartel screening
\citep{becker1968crime,stigler1970optimum,baker2002case,harrington2008detecting}.

The paper stops at the screening stage by design. Award-record
data does not contain what membership identification requires,
and cobidders are the validation object the data layer supports.
Adversarial adaptation degrades the screen's discrimination
predictably (Online Appendix~G), and replication rests on three
diagnostic conditions observable \emph{ex ante} in any candidate
jurisdiction: a winner/loser participation registry, an
electronic platform, and an institutional asymmetry between the
two sides of the bidding pool. The natural extensions are
cross-jurisdictional replication on platforms with comparable
observability layers, extension beyond commodity-heavy item
composition, and a formal characterisation of the optimal
sequencing of screening and forensic stages under endogenous
adversarial adaptation.

The conceptual content travels. Wherever an enforcement
environment exposes the award layer routinely while reserving
per-bidder bid amounts for forensic-recoverable access, the
right question is not whether bid-distribution screens can be
retrofitted onto data that does not include them at the layer
the screen must run on. The right question is what participation
primitive the collusive arrangement generates that survives in
the operational layer, and how the screening stage that reads
that primitive should be sequenced before the forensic stage
that requires the bid layer.
